\section{Prime representations and shifted Koszul cycles}
\label{sec:betti-euler-shifts}

Theorem~\ref{thm:fixed-weight-betti-euler} gives an algebraic estimate
for the linear Betti number in terms of ordered representation counts,
for any positive support.
Corollary~\ref{cor:prime-betti-singular-series} combines this estimate
with known prime asymptotics. Theorem~\ref{thm:prime-shift-proportion}
then constructs independent common-shift cycles spanning a positive
fraction of the linear strand at each sufficiently large weight of
the appropriate parity. We then deduce a Hilbert-series consequence
from the smaller Betti numbers in lower homological degrees.
For a positive support $\A$, let $u$ mark Goldbach degree and put
\[
 P(u)=\sum_{a\in\A}u^a,\qquad T_d(n)=[u^n]P(u)^d .
\]
Thus $T_d(n)$ counts \emph{ordered} representations of $n$ by $d$
labels in the specified support, with repetitions allowed.

\subsection{The fixed-weight Betti estimate}
Consider any $k$-ary Goldbach ring $G$ on a positive support $\A$.
Translation of labels as in Proposition~\ref{setup:affine-invariance}
reduces a bounded-below support to this case.

\begin{theorem}[Fixed-weight Betti estimate]
\label{thm:fixed-weight-betti-euler}
For every field and every $d\ge k\ge2$, put $i=d-k+1$. Then
\begin{equation}\label{eq:fixed-weight-leading}
 \beta_{i,n,d}
 =\frac{\binom{d-1}{i-1}}{d!}\,T_d(n)
      +O_d\bigl((1+|\A_n|)^{d-2}\bigr).
\end{equation}
\end{theorem}
\begin{proof}
Let $K=K(G)=\mathbb K_\kk\otimes_R G$ be the Koszul complex with
coefficients in $G$, and write $K_{i,n,d}$ for its term in homological
degree $i$ and bidegree $(n,d)$. Its differential is
$\mathrm d_{i,n,d}:K_{i,n,d}\to K_{i-1,n,d}$.
Since $\operatorname{im}\mathrm d_{i+1,n,d}\subseteq\ker\mathrm d_{i,n,d}$,
the rank--nullity theorem gives
\begin{equation}\label{eq:koszul-rank-nullity}
 \beta_{i,n,d}
 =\dim_\kk\bigl(K_{i,n,d}/\operatorname{im}\mathrm d_{i+1,n,d}\bigr)
  -\dim_\kk\operatorname{im}\mathrm d_{i,n,d}.
\end{equation}
We count the first term by grouping monomials modulo incoming
boundaries, then bound the second by the dimension of the target.

Recall that $i=d-k+1$ and $\deg\varepsilon_a=(a,1)$.
A term $f\varepsilon_\sigma$ in $K_{r,n,d}$, with $|\sigma|=r$,
therefore has coefficient bidegree
$\deg f=(n-\sum_{a\in\sigma}a,d-r)$. In the three relevant terms,
the coefficient monomial degrees are
\[
 d-(i+1)=k-2,\qquad d-i=k-1,\qquad d-(i-1)=k.
\]
Since $I$ is generated in monomial degree $k$, its components below
that degree vanish. Thus the coefficient spaces in $K_{i+1,n,d}$
and $K_{i,n,d}$ agree with those in the polynomial-ring Koszul complex.

For $x^\alpha\in\mathcal M_{n,d}$, put
$S=\operatorname{supp}\alpha$ and $s_\alpha=|S|$. For $\tau\subseteq S$,
write $u_\tau=(x^\alpha/x_\tau)\varepsilon_\tau$.
The $\binom{s_\alpha}{i}$ terms $u_\tau$ with $|\tau|=i$ form the
group belonging to $x^\alpha$. For
$\eta=\{a_1<\cdots<a_{i+1}\}\subseteq S$, the incoming column is
\[
 \mathrm d u_\eta
   =\sum_{j=1}^{i+1}(-1)^{j-1}u_{\eta\setminus\{a_j\}},
\]
the oriented simplex boundary on $S$.
To compute its rank, fix $a_0=\min S$ and take the boundaries
$\mathrm d u_{\{a_0\}\cup\tau}$ for the $i$-subsets
$\tau\subseteq S\setminus\{a_0\}$. These are independent:
each has coefficient $1$ at its own $u_\tau$, and all its other
exterior subsets contain $a_0$. They also span. If an $(i+1)$-subset $\eta$
omits $a_0$, the simplex identity $\partial^2=0$ on
$\{a_0\}\cup\eta$ expresses its boundary as a signed sum of
boundaries of faces containing $a_0$.
There are $\binom{s_\alpha-1}{i}$ selected boundaries, so this is
the incoming rank over every field. Consequently $x^\alpha$ contributes
\[
 \binom{s_\alpha}{i}-\binom{s_\alpha-1}{i}
   =\binom{s_\alpha-1}{i-1}
\]
dimensions modulo incoming boundaries. Summing gives
\[
 \dim_\kk(K_{i,n,d}/\operatorname{im}\mathrm{d}_{i+1,n,d})
   =\sum_{x^\alpha\in\mathcal M_{n,d}}
             \binom{s_\alpha-1}{i-1}.
\]
We take out-of-range binomial coefficients to be zero.

Let $T_d^{\mathrm{dist}}(n)$ count ordered representations of $n$
by $d$ pairwise distinct labels in $\A$, and put
$T_d^{\mathrm{rep}}(n)=T_d(n)-T_d^{\mathrm{dist}}(n)$.
Each squarefree monomial has exactly $d!$ orderings and has
$s_\alpha=d$. The combined squarefree contribution is therefore exactly
\[
 \frac{\binom{d-1}{i-1}}{d!}\,T_d^{\mathrm{dist}}(n).
\]
In the binary case $i=d-1$, each squarefree monomial has $d$ terms,
corresponding to the choice of a single coefficient label. One incoming boundary relates them,
leaving $d-1$ dimensions and the factor $(d-1)/d!$.

For a representation with repetitions, choose two equal positions
and fill the other $d-2$ positions with labels from $\A_n$.
The equation $2a+\sum_{\text{other positions}}a_j=n$ determines
the repeated label $a$ uniquely, if it exists. Hence
\[
 T_d^{\mathrm{rep}}(n)
 \le\binom d2(1+|\A_n|)^{d-2}.
\]
Every nonsquarefree monomial contributes a number of dimensions
bounded in terms of $d$, and their number is at most
$T_d^{\mathrm{rep}}(n)$. Thus both their contribution and the
correction from replacing $T_d^{\mathrm{dist}}(n)$ by $T_d(n)$
are absorbed in the stated error:
\[
 \begin{aligned}
 \dim_\kk(K_{i,n,d}/\operatorname{im}\mathrm{d}_{i+1,n,d})
 &=\frac{\binom{d-1}{i-1}}{d!}\,T_d^{\mathrm{dist}}(n)
       +O_d\bigl(T_d^{\mathrm{rep}}(n)\bigr)\\
 &=\frac{\binom{d-1}{i-1}}{d!}\,T_d(n)
       +O_d\bigl((1+|\A_n|)^{d-2}\bigr).
 \end{aligned}
\]

It remains to bound the second term in \eqref{eq:koszul-rank-nullity}.
In the target $K_{i-1,n,d}$, choosing the $i-1$ exterior labels fixes the residual
coefficient weight. Its monomial degree is $k$, and all represented
monomials of that degree and weight agree in $G$. Thus each exterior set gives
at most one target basis term, and
\[
 0\le\dim_\kk\operatorname{im}\mathrm d_{i,n,d}
 \le\dim_\kk K_{i-1,n,d}\le\binom{|\A_n|}{i-1}.
\]
Since $i-1=d-k\le d-2$, this loss fits the same error bound.
This proves \eqref{eq:fixed-weight-leading} for the Betti number.

\end{proof}

\subsection{Prime asymptotics and the cubic correction}

Theorem~\ref{thm:fixed-weight-betti-euler} uses only the Koszul complex
and elementary counting, and applies to finite as well as infinite
positive supports. For the full odd-prime support, the known
asymptotic for $T_d(n)$ makes its error relatively small when $d\ge3$
is fixed and $n\to\infty$ through $n\equiv d\pmod2$.
The following proof obtains the needed prime-counting asymptotic from
Tao and Ter\"av\"ainen's theorem \cite[Theorem~1.6]{TT}, then applies
the algebraic estimate to obtain the Betti asymptotic. The qualitative
linear-forms theorem comes from Green--Tao \cite{GT10,GT12} and
Green--Tao--Ziegler \cite{GTZ12}; the cited quantitative version
provides the uniform estimate used here and in the completion counts
below.

\begin{corollary}[Singular-series asymptotics]\label{cor:prime-betti-singular-series}
Let $G$ be the binary Goldbach ring on all odd primes, and fix $d\ge3$.
Define the $d$-prime singular series by
\begin{equation}\label{eq:d-prime-singular-series}
 \mathfrak S_d(n)=
 \prod_{\substack{s\text{ prime}\\s\mid n}}
       \left(1+\frac{(-1)^d}{(s-1)^{d-1}}\right)
 \prod_{\substack{s\text{ prime}\\s\nmid n}}
       \left(1-\frac{(-1)^d}{(s-1)^d}\right).
\end{equation}
As $n\to\infty$ through $n\equiv d\pmod2$, uniformly over coefficient fields,
\begin{equation}\label{eq:prime-betti-singular-asymptotic}
 \beta_{d-1,n,d}
 \sim\frac{\mathfrak S_d(n)}{d!(d-2)!}
             \frac{n^{d-1}}{(\log n)^d}.
\end{equation}
\end{corollary}
\begin{proof}
Apply the theorem of Tao and Ter\"av\"ainen
\cite[Theorem~1.6]{TT} to
$a_1,\ldots,a_{d-1},n-\sum_j a_j$ on the positive simplex
$a_j>0$, $\sum_j a_j<n$. Their linear parts are pairwise
nonproportional because $d\ge3$, and the affine constant $n$ is
allowed at scale $n$. At a prime $s$, the normalized local factor is
\[
 \frac{s}{(s-1)^d}
 \#\{(u_1,\ldots,u_d)\in(\mathbb F_s^\times)^d:
                                      u_1+\cdots+u_d=n\},
\]
which equals the corresponding factor in
\eqref{eq:d-prime-singular-series}. The factor at $2$ is $2$ on the
stated parity; the odd-prime factors give
$0<c_d\le\mathfrak S_d(n)\le C_d$ uniformly in $n$.
The simplex has volume $n^{d-1}/(d-1)!$, so the theorem yields
\[
 R_d(n):=\sum_{\substack{u_1+\cdots+u_d=n\\u_j\ge1}}
              \prod_{j=1}^d\Lambda(u_j)
 =\frac{\mathfrak S_d(n)}{(d-1)!}\,n^{d-1}+o_d(n^{d-1}).
\]
Tuples involving a proper prime power contribute at most
$O_d(n^{d-3/2}(\log n)^{d+1})=o(n^{d-1})$ to this weighted count.
Discard also tuples with an entry at most $n/(\log n)^{d+1}$.
Their unweighted number is $O_d(n^{d-1}/(\log n)^{d+1})$, and their
total weight is $O_d(n^{d-1}/\log n)$. Every remaining prime entry
is odd for large $n$ and has logarithm
$\log n+O_d(\log\log n)$. Hence
\begin{equation}\label{eq:d-prime-count-asymptotic}
 T_d(n)\sim\frac{\mathfrak S_d(n)}{(d-1)!}
                    \frac{n^{d-1}}{(\log n)^d}.
\end{equation}
Finally, $\pi(n)\ll n/\log n$ gives
\[
 \frac{(1+\pi(n))^{d-2}}{T_d(n)}
       \ll_d\frac{(\log n)^2}{n}\longrightarrow0.
\]
Theorem~\ref{thm:fixed-weight-betti-euler}, with coefficient
$(d-1)/d!$ for the binary ring, proves
\eqref{eq:prime-betti-singular-asymptotic}.
\end{proof}

For the binary ring on a positive support $\A$, at $d=3$ we can read
the correction directly from the cubic graph of
Example~\ref{ex:weight33}. Let $v_n$ and $c_n$ count its vertices and
components, where the vertices are the labels occurring in a triple
and $a\ne b$ are joined when $n-a-b\in\A$. A triple with three, two,
or one distinct labels contributes $2,1,0$ to
$\dim_\kk(K_{2,n,3}/\operatorname{im}\mathrm{d}_{3,n,3})$. Its outgoing
rank loss is the incidence rank $v_n-c_n$. Applying
\eqref{eq:koszul-rank-nullity} gives
\begin{equation}\label{eq:cubic-betti-euler}
 \beta_{2,n,3}
  =\frac{T_3(n)-\mathbf1_{\{\,n/3\in\A\,\}}}{3}-v_n+c_n.
\end{equation}
For example at weight $33$ there are $34$ ordered triples, including
$(11,11,11)$, and the graph is connected on eight vertices. The formula
gives $(34-1)/3-8+1=4$, as in the drawing.

\subsection{How many classes do common shifts give?}
Corollary~\ref{cor:prime-betti-singular-series} estimates the dimension
of the linear strand. We now identify explicit classes accounting for
a positive fraction of it. Theorem~\ref{thm:prime-shift-proportion}
selects an independent family from the common-shift cycles of
Section~\ref{sec:translations}, using estimates for simultaneous prime
patterns to count these classes.

Take $G$ to be the binary ring on the odd primes $\Podd$. Let
$W^+_{n,d}$ be the span in $\Tor^R_{d-1}(G,\kk)_{n,d}$ of the classes
$[z_{\sigma,h}]$ of Theorem~\ref{thm:translation}, for all $d$-element sets
$\sigma$ and all $h>0$ such that $\sigma,\sigma+h\subseteq\Podd$ and
$\sum_{a\in \sigma}a+h=n$. Different shifts can give dependent classes.
For example at weight $33$ the three possible positive shifts arise from
\[
 (\sigma,h)=(\{3,7,13\},10),\quad(\{3,7,19\},4),\quad(\{3,11,17\},2).
\]
These exhaust the possibilities by enumerating the prime triples of
weight $33$ in Example~\ref{ex:weight33}. In the first two cycles the
terms $x_{13}\varepsilon_7\varepsilon_{13}$ and
$x_7\varepsilon_7\varepsilon_{19}$ each occur only in their indicated
cycle, and no incoming Koszul boundary uses either term. The last
cycle is nonzero by Theorem~\ref{thm:translation}. Thus the shifts give
three of the four independent classes at this weight.

\begin{theorem}[Positive shifts at a fixed weight]
\label{thm:prime-shift-proportion}
For each $d\ge3$ there are constants $c_d,C_d>0$ depending only on $d$ such that, over every
field and for every sufficiently large $n\equiv d\pmod 2$,
\[
 c_d\,\frac{n^{d-1}}{(\log n)^d}
 \ \le\ \dim_\kk W^+_{n,d}\ \le\ \beta_{d-1,n,d}
 \ \le\ T_d(n)\ \le\ C_d\,\frac{n^{d-1}}{(\log n)^d}.
\]
In particular the positive shifts span at least a fixed positive fraction
of the linear strand in each such weight.
\end{theorem}
\begin{proof}
Fix $n\equiv d\pmod 2$; all intervals and completions below depend on this $n$.
Divide $[n/(4d),n/(3d)]$ into $d-1$ equal pieces, numbered from left
to right, and let $J_j$ be the middle half of the $j$th piece. Each
$J_j$ has length $n/(24d(d-1))$. We choose one odd prime $p_j\in J_j$
for $1\le j\le d-1$, so $\boldsymbol p=(p_1,\ldots,p_{d-1})$ is
strictly increasing. Set $J_0=[n/(8d),n/(6d)]$, to the left of $J_1$.

Call an odd prime $p_0\in J_0$ a \emph{completion} of $\boldsymbol p$ at weight $n$
if, with $h=n-\sum_{j=0}^{d-1}p_j$, every $p_j+h$ is prime
($0\le j\le d-1$). A completion gives the $d$ primes
$\sigma=\{p_0<p_1<\cdots<p_{d-1}\}$, and shifting them all by $h$ gives
the $d$ primes $\sigma+h$, so $z_{\sigma,h}$ is a translation cycle of
weight $n$. The smallest shifted prime is
\[
 p_0+h=n-\sum_{j=1}^{d-1}p_j .
\]
Thus, with $n$ fixed, $p_0+h$ depends only on $\boldsymbol p$, and it exceeds $2n/3$.
For each tuple admitting a completion, choose its least such $p_0$.
This choice determines $h$; the shift is not chosen independently and
may vary with the tuple.

The selected cycles are linearly independent, although different
tuples may have different shifts.
In the cycle for $\boldsymbol p$, consider the term
$x_{p_0+h}\varepsilon_{p_1}\cdots\varepsilon_{p_{d-1}}$. Every other term
of this cycle has $p_0\in J_0$ in its exterior part. Thus each selected
cycle has just one term whose exterior labels all lie in
$J_1,\ldots,J_{d-1}$, and those labels are the entries of $\boldsymbol p$.
Two tuples with the same sum share the coefficient label $p_0+h$, but
their distinguished terms still differ in the exterior part.
The least-completion rule then determines $p_0$, $h$, and the selected cycle.
The sole possible incoming Koszul column
containing it is $\mathrm{d}(\varepsilon_\tau)$ for
$\tau=\{p_1<\cdots<p_{d-1}<p_0+h\}$, since the incoming coefficients have monomial degree zero. The same column has a term with $p_0+h$ in
the exterior part. No selected cycle has such a term, and the
coefficient label and exterior indices determine this incoming column
uniquely. Thus in any boundary equal to a combination of the selected
cycles the coefficient of this column must vanish. Reading the coefficient
of the distinguished term proves the claimed independence over every field.

Let $\mathcal T_+$ be the set of tuples $\boldsymbol p$ admitting
a completion, and let $f(\boldsymbol p)$ count all their completions.
Write
\[
 \langle f,1\rangle=\sum_{\boldsymbol p\in\mathcal T_+}f(\boldsymbol p),
 \qquad
 \langle f,f\rangle=\sum_{\boldsymbol p\in\mathcal T_+}f(\boldsymbol p)^2.
\]
These count completions and ordered pairs of completions of a common
tuple, respectively, including repeated pairs. By Cauchy--Schwarz,
$|\mathcal T_+|\ge\langle f,1\rangle^2/\langle f,f\rangle$.
We will prove
\begin{equation}\label{eq:shift-moments}
 \langle f,1\rangle\gg_d\frac{n^d}{(\log n)^{2d}},\qquad
 \langle f,f\rangle\ll_d\frac{n^{d+1}}{(\log n)^{3d}}.
\end{equation}
The selected independent classes then give
\[
 \dim_\kk W^+_{n,d}\ge|\mathcal T_+|
 \ge\frac{\langle f,1\rangle^2}{\langle f,f\rangle}
 \gg_d\frac{n^{d-1}}{(\log n)^d}.
\]

Both moments count simultaneous prime values of linear forms on a
box. We use Tao and Ter\"av\"ainen \cite[Thm.~1.6]{TT}: for bounded,
pairwise nonproportional slopes and constant terms $O(n)$, the
von Mangoldt weighted count has a uniform asymptotic whose main term
is the box volume times the product of local factors.
With $D$ free variables and $t$ forms, its error here is
$O_d(n^D(\log\log n)^{-\eta_d})$ for some $\eta_d>0$.
We use $(D,t)=(d,2d)$ and $(d+1,3d)$.
Since all forms on our boxes are comparable to $n$, converting weighted
bounds to prime counts divides both the main-term and error scales by
$(\log n)^t$. Thus the saving in $\log\log n$ suffices.

\emph{The first bound.} Use one integer variable $a_j\in J_j$ for each
$0\le j\le d-1$; the index $j$ specifies the interval in which the variable
ranges. A point $\boldsymbol a=(a_0,\ldots,a_{d-1})$ of
$J_0\times\cdots\times J_{d-1}$ gives a tuple and a completion, with
$p_j=a_j$, precisely when all $2d$ forms
\[
 a_j,\qquad n-\sum_{\substack{0\le r\le d-1\\r\ne j}}a_r
 \qquad (0\le j\le d-1)
\]
take prime values; the second family is $a_j+h$. Thus $\langle f,1\rangle$
is the number of these points.
Their slopes are pairwise nonproportional because $d\ge3$.
At an odd prime $s$ choose nonzero residues for the $a_j$ whose sum
is $n$ modulo $s$. This is possible since every residue is the sum
of two nonzero residues. The shifted forms then also have nonzero
residues. At $2$ take all $a_j$ odd and use $n\equiv d\pmod 2$.
These choices bound the finitely many small local factors below
uniformly in $n$. At all sufficiently large primes $s$ the slopes remain
pairwise independent modulo $s$. Counting intersections of pairs
of their zero hyperplanes gives normalized local factors
$1+O_d(s^{-2})$. Hence the product of the local factors is uniformly
bounded below. Our box has volume $\gg_d n^d$, and all forms there
are comparable to $n$. The contribution of proper prime powers is
$o(n^d)$ in the weighted count: one form takes at most
$O(\sqrt n\log n)$ such values, each on $O_d(n^{d-1})$ points.
Including the $2d$ logarithmic weights gives a total of at most
$O_d(n^{d-1/2}(\log n)^{2d+1})=o_d(n^d)$.
Converting logarithmic weights to an unweighted count gives the first
bound in \eqref{eq:shift-moments}.

\emph{The second bound.} A pair of completions $(p_0,p'_0)$ of a common
tuple is a point of $J_0\times J_0\times J_1\times\cdots\times J_{d-1}$:
two integer variables $a_0,a'_0\in J_0$ and the same $a_j\in J_j$ for
$1\le j\le d-1$. Put $a_d=n-\sum_{j=1}^{d-1}a_j$, which takes the value
$p_0+h$ at $\boldsymbol p$. The point is counted by $\langle f,f\rangle$
precisely when the forms
\[
 a_d,\ a_0,\ a'_0,\quad
 a_j,\ a_d-a_0+a_j,\ a_d-a'_0+a_j \qquad(1\le j\le d-1)
\]
are prime. There are $d+1$ free variables because the two completions share the
$d-1$ exterior labels. Their prime requirements also share those labels
and $a_d$, leaving $4d-d=3d$ distinct forms. Their slopes are pairwise
nonproportional. Each of the last two families uses its own completion
variable, and the form indexed by $j$ omits precisely $a_j$.
Coincident completions $a_0=a'_0$ are included: equality of some values
at such points does not change these linear parts, and the weighted
product still contains all $3d$ logarithmic factors. The same pairwise intersection
estimate bounds the product of the local factors above uniformly.
Each prime value in this box has logarithmic weight comparable to
$\log n$, so the linear forms theorem gives the second bound in
\eqref{eq:shift-moments}.

The entire Koszul module in homological degree $d-1$ has at most $T_d(n)$
basis terms, since placing the coefficient label first and the
increasing exterior labels after it gives an ordered representation.
The upper bound $T_d(n)\ll_d n^{d-1}/(\log n)^d$ follows from
\eqref{eq:d-prime-count-asymptotic} and the uniform upper bound for
$\mathfrak S_d(n)$.
\end{proof}

\subsection{A Hilbert-series consequence}
\label{sec:hilbert-euler-consequence}
The preceding estimates concern individual Betti numbers and explicit
cycles. Their alternating sum connects this growth with the Hilbert
series. For a $k$-ary Goldbach ring on a positive support, define
\[
 \gamma_{n,d}=\sum_{i\ge0}(-1)^i\beta_{i,n,d},\qquad
 E_G(u,v)=\sum_{n,d}\gamma_{n,d}u^nv^d.
\]
Let $H_G(u,v)=\sum_{n,d}(\dim_\kk G_{n,d})u^nv^d$ and let $H_R$
be the corresponding series of $R$. Taking the Euler characteristic
of the Koszul complex in each bidegree gives
\begin{equation}\label{eq:betti-euler-hilbert}
 E_G(u,v)=\frac{H_G(u,v)}{H_R(u,v)}
        =H_G(u,v)\prod_{a\in\A}(1-u^av).
\end{equation}
Each coefficient is a finite sum. By the component basis in
Proposition~\ref{setup:components}, $H_G$ is independent of the
coefficient field, and hence so are $E_G$ and $\gamma_{n,d}$.

\begin{corollary}[Comparison with the Euler coefficient]
\label{cor:hilbert-euler-comparison}
For every field and every $d\ge k\ge2$, put $i=d-k+1$. Then
\[
 \bigl|(-1)^i\gamma_{n,d}-\beta_{i,n,d}\bigr|
 \le\sum_{j<i}\beta_{j,n,d}
 =O_d\bigl((1+|\A_n|)^{i-1}\bigr).
\]
Thus $(-1)^i\gamma_{n,d}$ has the same leading term as
$\beta_{i,n,d}$ in Theorem~\ref{thm:fixed-weight-betti-euler}.
For the binary ring on all odd primes, fixed $d\ge3$, and
$n\to\infty$ through $n\equiv d\pmod2$,
\[
 (-1)^{d-1}\gamma_{n,d}\sim\beta_{d-1,n,d}
 \sim\frac{d-1}{d!}T_d(n).
\]
\end{corollary}
\begin{proof}
Write $K=K(G)$ for the Koszul complex, with the conventions of
Theorem~\ref{thm:fixed-weight-betti-euler}.
In homological degrees $j>i$ the Koszul differential and both adjacent terms
in this bidegree agree with the exact polynomial-ring complex in
positive homological degrees, so $\beta_{j,n,d}=0$. Therefore
\[
 (-1)^i\gamma_{n,d}=\beta_{i,n,d}
       +\sum_{j<i}(-1)^{i+j}\beta_{j,n,d}.
\]
For $j<i$ the coefficients have monomial degree
$d-j\ge k$. Choose the $j$ exterior labels and $d-j-k$ coefficient
labels. The remaining $k$ labels have prescribed sum and contribute
at most one coefficient class. These choices span $K_{j,n,d}$ and require
choosing only $j+(d-j-k)=d-k=i-1$ labels. Hence
\[
 \beta_{j,n,d}\le\dim_\kk K_{j,n,d}
   =O_d\bigl((1+|\A_n|)^{i-1}\bigr).
\]
There are only finitely many such homological degrees, bounded in number by $d$.
It follows that
\[
 |(-1)^i\gamma_{n,d}-\beta_{i,n,d}|
 \le\sum_{j<i}\beta_{j,n,d}
 =O_d\bigl((1+|\A_n|)^{i-1}\bigr).
\]
Since $i-1=d-k\le d-2$, Theorem~\ref{thm:fixed-weight-betti-euler}
gives the asserted leading term. For binary odd-prime support,
Corollary~\ref{cor:prime-betti-singular-series} makes the error
negligible relative to that term.
\end{proof}

For a binary ring, the ideal is generated in monomial degree two
and $\Tor_0^R(G,\kk)$ has monomial degree zero. Consequently
$\gamma_{n,3}=\beta_{2,n,3}$, so \eqref{eq:cubic-betti-euler}
also computes this Euler coefficient exactly. In monomial degree four,
$-\gamma_{n,4}=\beta_{3,n,4}-\beta_{2,n,4}$.
The identity \eqref{eq:betti-euler-hilbert} therefore recovers
the dominant Betti growth from Hilbert-series coefficients, while
retaining the cancellations between homological degrees.
The same Euler-characteristic principle compares binary and ternary
rings through mixed Tor in \cite{CompanionTopology}.
